% =========================================================================
\section{Introducción}\label{sec:introduccion}
% =========================================================================
El electromagnetismo clásico es, en esencia, una teoría de \emph{campos}: a
cada punto del espacio se le asigna una magnitud (un número, un vector) que
varía de un punto a otro y, en general, también con el tiempo. El campo
eléctrico $\vc{E}(\vc{r})$ y el campo magnético $\vc{B}(\vc{r})$ son campos
vectoriales; el potencial electrostático $V(\vc{r})$ y la densidad de carga
$\rho(\vc{r})$ son campos escalares. Para hablar de cómo cambian esos campos
(\emph{derivadas}), de cuánto «atraviesan» una superficie o «circulan» a lo
largo de un circuito (\emph{integrales}), y de cómo conectar ambas ideas,
necesitamos un pequeño conjunto de herramientas matemáticas, que son las que
reúne este tema~\cite{NAVAS,GRIFFITHS,ARFKEN}. Las figuras reproducen, ahora como
diagramas vectoriales en TikZ y gráficas en pgfplots, las del material de
figuras del curso~\cite{FIGURAS-T0}, que se prepararon con Mathematica, y las completan con
diagramas propios de los ejercicios de verificación.

Estas herramientas se pueden agrupar en tres bloques que se apoyan unos en
otros:
\begin{itemize}
  \item \textbf{Cálculo diferencial vectorial} (\cref{sec:campos,sec:operadores}):
        el gradiente de un campo escalar, y la divergencia y el rotacional de un
        campo vectorial, con sus reglas de derivación y las identidades que
        satisfacen entre sí.
  \item \textbf{Cálculo integral vectorial} (\cref{sec:integral}): las integrales
        de línea y de flujo, y los tres teoremas integrales fundamentales
        (del gradiente, de Stokes y de Gauss), todos ellos generalizaciones del
        teorema fundamental del cálculo.
  \item \textbf{Geometría de las coordenadas} (\cref{sec:coordenadas,sec:posicion}):
        los sistemas cartesiano, cilíndrico y esférico con sus operadores
        diferenciales, y las funciones de la posición relativa
        $\vc{R}=\vc{r}-\vc{r}\,'$ que aparecen siempre que una distribución
        de fuentes genera un campo, junto con la delta de Dirac.
\end{itemize}
El documento termina con unos \emph{seminarios prácticos} (\cref{sec:seminarios})
donde se resuelven, con todo detalle, los ejercicios de verificación de los
teoremas y las demostraciones de identidades vectoriales con el tensor de
Levi-Civita. Los resultados se enuncian con el rigor de un artículo, pero
procuramos acompañar cada fórmula de una explicación geométrica o física que
permita \emph{ver} por qué es cierta.

% -------------------------------------------------------------------------
\subsection{Notación y convenios}\label{sec:notacion}
% -------------------------------------------------------------------------
Seguimos la notación vectorial de L.~Navas~\cite{NAVAS} y las expresiones del
formulario del curso~\cite{FORMULARIO}. Las flechas distinguen los
vectores de los escalares, y los versores (vectores unitarios) se escriben con
sombrero.

\begin{table}[ht]
\centering
\caption{Resumen de la notación empleada.}
\label{tab:notacion}
\small
\renewcommand{\arraystretch}{1.35}
\begin{tabular}{@{}p{4.6cm}p{9.2cm}@{}}
\toprule
Símbolo & Significado\\
\midrule
$\vc{r}=x\,\ux+y\,\uy+z\,\uz$ & Vector de posición del \emph{punto campo};
  $\vc{r}\,'$ es el del \emph{punto fuente}.\\
$\ux,\uy,\uz$;\ $\ur,\uth,\uph$ & Versores de las bases cartesiana y esférica
  (en la base cilíndrica, $\urho,\uph,\uz$).\\
$f,\,g,\,V$ & Campos escalares, $f:W\subseteq\R^{n}\to\R$.\\
$\vc{F},\,\vc{A},\,\vc{B}$ & Campos vectoriales, $\vc{F}:W\subseteq\R^{n}\to\R^{n}$,
  con componentes $(F_x,F_y,F_z)$.\\
$\nab=\ux\,\partial_x+\uy\,\partial_y+\uz\,\partial_z$ & Operador «nabla»
  (o «del»), un operador diferencial; $\nab'$ actúa sobre las coordenadas fuente.\\
$\nab f,\ \nab\cdot\vc{F},\ \nab\times\vc{F},\ \Lap f$ & Gradiente, divergencia,
  rotacional y laplaciano.\\
$\dl,\ \dS,\ \dV$ & Elemento de línea dirigido, de área dirigido
  ($\dS=\un\,\dsc$) y de volumen.\\
$\partial S,\ \partial V$ & Borde (orientado) de una superficie $S$ y de un
  volumen $V$.\\
$\displaystyle\int_{C}\!,\ \oint_{C},\ \iint_{S},\ \oiint_{S},\ \iiint_{V}$ &
  Integrales sobre curvas abiertas y cerradas, sobre superficies abiertas y
  cerradas, y sobre volúmenes.\\
$\varepsilon_{ijk},\ \delta_{ij}$ & Tensor de Levi-Civita y delta de Kronecker.\\
$\delta^{3}(\vc{r})$ & Delta de Dirac tridimensional.\\
\bottomrule
\end{tabular}
\end{table}

\begin{notacion}[Convenios generales]\label{not:convenios}
Fijamos los siguientes convenios.\begin{enumerate}[label=(\alph*)]
  \item Salvo mención en contrario, los campos son de clase $C^{1}$ en la región
        considerada (de clase $C^{2}$ cuando intervienen derivadas segundas), y
        las curvas y superficies son regulares a trozos.
  \item Las \emph{regiones} $W\subseteq\R^{n}$ son abiertos conexos.
  \item Las curvas y las superficies están \emph{orientadas}. El borde de una
        superficie lleva la orientación compatible con la de la superficie
        según la regla de la mano derecha (\cref{sec:stokes}), y el borde de
        un sólido se orienta con el vector normal exterior.
  \item En los desarrollos con índices se aplica el convenio de sumación de
        Einstein: todo índice repetido en un mismo monomio se suma de $1$ a $3$,
        y $\partial_i\equiv\partial/\partial x_i$ con
        $(x_1,x_2,x_3)=(x,y,z)$.
  \item Escribimos $\sen$ para el seno y $\tan^{-1}$ para la función arcotangente.
\end{enumerate}
\end{notacion}

\begin{observacion}[Sobre el símbolo $\nab$]\label{obs:nabla}
El símbolo $\nab$ se manipula como si fuese un vector de componentes
$(\partial_x,\partial_y,\partial_z)$, y esa analogía es muy útil para recordar
fórmulas: $\nab f$ «es» el vector $\nab$ multiplicado por el escalar $f$,
$\nab\cdot\vc{F}$ su producto escalar con $\vc{F}$, y $\nab\times\vc{F}$ su
producto vectorial. Pero conviene no olvidar que $\nab$ \emph{no es un vector}
hasta que opera sobre un campo: es un operador diferencial, y en particular no
conmuta con los campos ($\nab\cdot\vc{A}\neq\vc{A}\cdot\nab$). La analogía
sugiere fórmulas, pero no las demuestra.
\end{observacion}
